\documentclass[a4paper,10pt]{article}
\usepackage[utf8]{inputenc}
\usepackage[T1]{fontenc}
\usepackage[french]{babel}
\usepackage[margin=1.5cm]{geometry}
\usepackage{xcolor}
\usepackage{tcolorbox}
\usepackage{enumitem}
\usepackage{multicol}
\usepackage{lmodern}
\usepackage{tabularx}
\usepackage{booktabs}

\renewcommand{\familydefault}{\sfdefault}
\pagestyle{empty}
\setlist{nosep, leftmargin=1.5em}

\definecolor{navy}{RGB}{0,51,102}
\definecolor{teal}{RGB}{0,120,120}
\definecolor{crimson}{RGB}{180,0,0}
\definecolor{green}{RGB}{0,120,50}
\definecolor{lightbg}{RGB}{245,248,252}
\definecolor{orange}{RGB}{200,100,0}
\definecolor{purple}{RGB}{100,0,160}

\tcbset{basebox/.style={
  colback=lightbg, colframe=navy, coltitle=white,
  fonttitle=\bfseries, arc=2mm, boxrule=0.5mm,
  left=3mm, right=3mm, top=2mm, bottom=2mm, after skip=6pt
}}

\newcommand{\key}[1]{\textbf{\textcolor{navy}{#1}}}
\newcommand{\warn}[1]{\textbf{\textcolor{crimson}{#1}}}
\newcommand{\ok}[1]{\textbf{\textcolor{green}{#1}}}

\begin{document}

\begin{center}
  \colorbox{green}{\parbox{\dimexpr\textwidth-2\fboxsep}{
    \centering\vspace{3pt}
    {\LARGE\bfseries\textcolor{white}{GESTION DU SI — Patrimoine, Incidents, GLPI}}\\
    {\normalsize\textcolor{white}{BTS SIO — Maël Mignard}}
    \vspace{3pt}
  }}
\end{center}
\vspace{6pt}

\begin{multicols}{2}

% ============================================================
\begin{tcolorbox}[basebox, colframe=green, title=1. Gestion du Patrimoine Informatique]
\key{Patrimoine informatique} = ensemble des ressources matérielles et logicielles d'une organisation.\\[4pt]
\textbf{Composantes :}
\begin{itemize}
  \item \key{Matériel} : PC, serveurs, imprimantes, routeurs, switches
  \item \key{Logiciel} : OS, applications, licences
  \item \key{Réseau} : câblage, équipements actifs
  \item \key{Données} : fichiers, bases de données, sauvegardes
\end{itemize}
\vspace{4pt}
\textbf{Cycle de vie d'un équipement :}
\begin{enumerate}
  \item Achat / Approvisionnement
  \item Réception et inventaire
  \item Déploiement / Installation
  \item Utilisation et maintenance
  \item Mise hors service / Recyclage
\end{enumerate}
\vspace{4pt}
\key{Inventaire} : recensement exhaustif du parc \\
\key{Numéro de série} : identification unique du matériel
\end{tcolorbox}

% ============================================================
\begin{tcolorbox}[basebox, colframe=green, title=2. GLPI — Présentation]
\key{GLPI} (Gestionnaire Libre de Parc Informatique) :\\
Outil open-source de gestion du parc et des tickets.\\[4pt]
\textbf{Fonctionnalités principales :}
\begin{itemize}
  \item \ok{Inventaire matériel et logiciel}
  \item \ok{Gestion des tickets (incidents et demandes)}
  \item \ok{Base de connaissances}
  \item \ok{Gestion des contrats et fournisseurs}
  \item \ok{Statistiques et tableaux de bord}
\end{itemize}
\vspace{4pt}
\key{FusionInventory} : plugin d'inventaire automatique \\
\key{Agent GLPI} : logiciel installé sur chaque poste \\
\quad $\to$ remonte les informations matérielles/logicielles \\[4pt]
\key{CMDB} (Configuration Management Database) : base centralisant toutes les configurations
\end{tcolorbox}

% ============================================================
\begin{tcolorbox}[basebox, colframe=teal, title=3. GLPI — Gestion des tickets]
\textbf{Types de tickets :}
\begin{itemize}
  \item \key{Incident} : dysfonctionnement signalé
  \item \key{Demande} : requête de service (installation, accès...)
\end{itemize}
\vspace{4pt}
\textbf{Cycle de vie d'un ticket :}
\begin{enumerate}
  \item Ouverture (création par l'utilisateur ou technicien)
  \item Attribution (à un technicien ou groupe)
  \item Traitement (diagnostic, résolution)
  \item Validation / Clôture
\end{enumerate}
\vspace{4pt}
\textbf{Informations d'un ticket :}
\begin{itemize}
  \item Titre et description
  \item Priorité (très haute / haute / moyenne / basse)
  \item Urgence et Impact
  \item Catégorie
  \item Technicien assigné
  \item Date d'ouverture / clôture
\end{itemize}
\vspace{4pt}
\key{SLA} (Service Level Agreement) : délais de traitement contractuels
\end{tcolorbox}

% ============================================================
\begin{tcolorbox}[basebox, colframe=crimson, title=4. Gestion des Incidents — ITIL]
\key{ITIL} (Information Technology Infrastructure Library) :\\
Référentiel de bonnes pratiques pour la gestion IT.\\[4pt]
\textbf{Différencier les situations :}
\begin{tabular}{ll}
\toprule
\textbf{Type} & \textbf{Définition} \\
\midrule
\key{Incident} & Interruption non planifiée d'un service \\
\key{Problème} & Cause inconnue d'un ou plusieurs incidents \\
\key{Demande} & Requête standard de service \\
\key{Changement} & Modification planifiée et contrôlée \\
\bottomrule
\end{tabular}
\vspace{4pt}\\
\textbf{Exemples :}
\begin{itemize}
  \item \warn{Incident} : PC qui ne démarre plus
  \item \warn{Problème} : panne réseau répétée chaque lundi
  \item \ok{Demande} : installer un logiciel
  \item \ok{Changement} : migrer un serveur
\end{itemize}
\end{tcolorbox}

% ============================================================
\begin{tcolorbox}[basebox, colframe=crimson, title=5. Analyse d'impact et priorité]
\key{Impact} : étendue des perturbations causées
\begin{itemize}
  \item \warn{Critique} : arrêt total de l'activité (serveur central)
  \item Élevé : service majeur impacté (messagerie HS)
  \item Moyen : contournement possible (imprimante HS)
  \item Faible : utilisateur unique impacté
\end{itemize}
\vspace{4pt}
\key{Urgence} : vitesse à laquelle résoudre\\[4pt]
\key{Priorité} = f(Impact, Urgence)\\[4pt]
\textbf{Communication en situation de crise :}
\begin{itemize}
  \item Garder son calme et faire preuve d'empathie
  \item Reformuler pour montrer la compréhension
  \item Rester neutre, courtois et rassurant
  \item Éviter le jargon technique avec l'utilisateur
  \item Ne jamais rejeter la faute sur l'utilisateur
\end{itemize}
\end{tcolorbox}

% ============================================================
\begin{tcolorbox}[basebox, colframe=navy, title=6. Processus de résolution d'incident]
\begin{enumerate}
  \item \key{Identification} : détecter et signaler l'incident
  \item \key{Enregistrement} : créer un ticket avec toutes les infos
  \item \key{Classification} : catégorie, priorité, impact
  \item \key{Diagnostic initial} : first-level (hotline)
  \item \key{Escalade} si nécessaire vers niveau 2 ou 3
  \item \key{Résolution} : corriger le problème
  \item \key{Clôture} : confirmer avec l'utilisateur
  \item \key{Retour d'expérience} : documenter la solution
\end{enumerate}
\vspace{4pt}
\key{Niveaux de support :}
\begin{itemize}
  \item N1 : Hotline, 1er contact (réinitialisation MDP...)
  \item N2 : Technicien spécialisé (installation, config)
  \item N3 : Expert/éditeur (bug logiciel, infrastructure)
\end{itemize}
\end{tcolorbox}

% ============================================================
\begin{tcolorbox}[basebox, colframe=orange, title=7. Gestion des licences logicielles]
\key{Licence propriétaire} : coût, restrictions d'usage \\
\key{Logiciel libre} (open-source) : code accessible, modifiable \\[4pt]
\textbf{Types de licences :}
\begin{itemize}
  \item \key{OEM} : liée au matériel (ne peut pas être transférée)
  \item \key{Monoposte} : 1 installation sur 1 PC
  \item \key{Volume} : entreprise, nombre de postes
  \item \key{SaaS} : abonnement, logiciel en ligne
  \item \key{GPL / MIT / Apache} : licences open-source
\end{itemize}
\vspace{4pt}
\warn{Utiliser un logiciel sans licence = délit} \\
\key{Audit logiciel} : vérifier la conformité des licences \\
\key{SAM} (Software Asset Management) : gestion des actifs logiciels
\end{tcolorbox}

% ============================================================
\begin{tcolorbox}[basebox, colframe=purple, title=8. Masterisation et déploiement]
\key{Masterisation} : créer une image système de référence \\
\quad $\to$ Cloner sur de nombreux postes rapidement \\[4pt]
\textbf{Outils :}
\begin{itemize}
  \item \key{Clonezilla} : clonage disque open-source
  \item \key{FOG} : déploiement réseau (PXE boot)
  \item \key{MDT / WDS} : déploiement Windows
  \item \key{Ansible / Puppet} : automatisation configuration
\end{itemize}
\vspace{4pt}
\key{PXE} (Preboot Execution Environment) : démarrage réseau \\
\key{Sysprep} : préparer Windows pour le clonage \\
\key{WSUS} : serveur de mises à jour Windows en entreprise
\end{tcolorbox}

% ============================================================
\begin{tcolorbox}[basebox, colframe=teal, title=9. Télémaintenance et assistance]
\key{Télémaintenance} : intervention à distance sur un poste \\[4pt]
\textbf{Outils de prise en main à distance :}
\begin{itemize}
  \item \key{RDP} (Remote Desktop Protocol) : Windows, port 3389
  \item \key{VNC} : multiplateforme
  \item \key{TeamViewer / AnyDesk} : solutions commerciales
  \item \key{SSH} : ligne de commande Linux, port 22
\end{itemize}
\vspace{4pt}
\textbf{Procédure :}
\begin{enumerate}
  \item Identifier l'utilisateur et son problème
  \item Obtenir son accord pour la prise en main
  \item Intervenir en expliquant les actions
  \item Clore la session et documenter
\end{enumerate}
\warn{Toujours demander l'autorisation avant de prendre la main}
\end{tcolorbox}

% ============================================================
\begin{tcolorbox}[basebox, colframe=green, title=10. Documentation et base de connaissances]
\key{Base de connaissances} : recueil de solutions aux incidents récurrents \\[4pt]
\textbf{Avantages :}
\begin{itemize}
  \item Résolution plus rapide (solution déjà documentée)
  \item Autonomie des utilisateurs (self-service)
  \item Capitalisation des compétences
\end{itemize}
\vspace{4pt}
\textbf{Format d'une fiche de résolution :}
\begin{itemize}
  \item Titre (symptôme)
  \item Catégorie / Environnement
  \item Description du problème
  \item Cause identifiée
  \item Solution étape par étape
  \item Date et auteur
\end{itemize}
\vspace{4pt}
\key{GLPI} intègre une base de connaissances native \\
\key{Confluence / Notion} : outils de documentation wiki
\end{tcolorbox}

\end{multicols}

\vspace{4pt}
\begin{tcolorbox}[basebox, colframe=green, title=Points clés à retenir pour l'examen]
\begin{multicols}{2}
\begin{itemize}
  \item \key{GLPI} = gestion de parc + ticketing, open-source
  \item \key{ITIL} = référentiel de bonnes pratiques IT
  \item Incident ≠ Problème ≠ Demande ≠ Changement
  \item Incident = interruption non planifiée ; Problème = cause inconnue récurrente
  \item Priorité = f(Impact, Urgence)
  \item N1/N2/N3 : niveaux d'escalade du support
  \item \key{SLA} = engagement contractuel sur les délais de résolution
  \item Cycle de vie : Achat $\to$ Déploiement $\to$ Maintenance $\to$ Fin de vie
  \item Licence OEM = liée au matériel, non transférable
  \item Masterisation = image disque de référence pour clonage de masse
  \item RDP port 3389 = accès bureau à distance Windows
  \item \warn{Jamais de prise en main sans accord de l'utilisateur}
\end{itemize}
\end{multicols}
\end{tcolorbox}

\end{document}
